%%%%%%%%%%%%%%%%%%%%%%%%%%%%%%%%%
%
%       EXERCÍCIO
%
%%%%%%%%%%%%%%%%%%%%%%%%%%%%%%%%%

\ifdefstring{\atividade}{prova}{%
    \ifdefstring{\modo}{objetivo}{%
        \renewcommand{\valorquestao}{\ValorQObj\ ponto}
    }{%
        \renewcommand{\valorquestao}{\ValorQDisc\ pontos}
    }%
}%

\begin{exercicioBanco}[\valorquestao]
Sejam \(A = \{x \in \bbR \mid -5 \le x \le 2\}\) e \(B = \{x \in \bbR \mid 0 \le x < 4\}\). Determine \(A\setminus B\).

% Define as alternativas
\newcommand{\alternativas}{%
    \begin{center}
        \begin{tabularx}{\textwidth}{XXX}
            (a) \(\{x \in \bbR \mid 2 < x < 4\}\). &
            (b) \(\{x \in \bbR \mid 0 \le x \le 2\}\). &
            (c) \(\{x \in \bbR \mid -5 \le x < 4\}\). \\[5pt]
            (d) \(\{x \in \bbR \mid -5 \le x < 0\}\). &
            \multicolumn{2}{l}{(e) \(\{x \in \bbR \mid -5 \le x < 0 \; \text{ ou } \; 2 < x < 4\}\).}
        \end{tabularx}
    \end{center}
}

% Define a resposta correta
\newcommand{\resposta}{D}

% Lógica condicional para exibição
\ifdefstring{\atividade}{lista}{%
    \alternativas
    \vspace{0.5em}
    
    \noindent\textbf{Resposta:} letra \textbf{\resposta}.
}{%
    \ifdefstring{\modo}{objetivo}{%
        \alternativas
    }{%
        % modo = discursiva → não mostra alternativas
    }
}
\end{exercicioBanco}

